\documentclass[12pt]{article}
\usepackage[utf8]{inputenc}
\usepackage[T1]{fontenc}
\usepackage{hyperref}
\title{Wakawaka Language Resources for K-12 Curriculum Development}
\author{Compiled by Claude Sonnet 5.5 (claude-sonnet-5-5) for Andr\'as Kornai}
\date{}
\begin{document}
\maketitle

\section{Introduction}
Wakawaka (Wakka Wakka; ISO 639-3 wkw; Glottocode waka1274) is a Pama-Nyungan language of south-eastern Queensland, Australia, with a country around the Burnett region. It is dormant or extinct as a first language. Wikipedia gives 1965 as the date of the last fluent speaker's death, and the name Duungidjawu is listed as a dialect. Documentation is mainly historical wordlists and grammatical sketches, plus one modern grammar of the Duungidjawu variety. This report is based on a single round of web searches. It lists only what those searches returned, and I opened none of the linked pages.

\section{Primary Educational Resources}
No published primer or textbook for Wakawaka turned up in the searches. A search result says Eidsvold State School has introduced lessons for primary-age students, but I found no curriculum documents. A revival effort is mentioned on the Wikipedia page: \url{https://en.wikipedia.org/wiki/Wakka_Wakka_language}.

\section{Linguistic and Reference Materials}
\begin{itemize}
\item State Library of Queensland, ``Wakka Wakka Body Parts'' word list (PDF): \url{https://content.slq.qld.gov.au/sites/default/files/Wakka%20Wakka%20Body%20Parts.pdf}
\item State Library of Queensland, Aboriginal and Torres Strait Islander languages word lists (data.gov.au catalogue entry): \url{https://data.gov.au/data/dataset/3fa7a159-f280-48f9-9a83-9209f73942ef}
\item University of Queensland Fryer Library manuscript catalogue entry for the Wakawaka language, which includes an English--Wakka and Wakka--English vocabulary in Watson's papers: \url{https://manuscripts.library.uq.edu.au/index.php/wakawaka-language}
\item University of Queensland manuscripts catalogue, language vocabulary word lists: \url{https://manuscripts.library.uq.edu.au/index.php/language-vocabulary-word-lists}
\item WorldCat record found in the search (not checked): \url{https://search.worldcat.org/oclc/930955155}
\end{itemize}
Other works named in the search summaries, for which I saw no direct URL: Kite and Wurm (2004), \emph{The Duungidjawu Language of Southeast Queensland: Grammar, Texts and Vocabulary}; Holmer (1983), \emph{Linguistic Survey of South-Eastern Queensland}; Watson's \emph{Four vocabularies of four representative tribes of South-Eastern Queensland}; and the grammatical work of Rev. John Mathew.

\section{Religious and Community Resources}
I found no scripture translations or community organisation pages for Wakawaka in these searches.

\section{Digital and Audio Resources}
I found no audio or video resources and no software for Wakawaka in these searches.

\section{Curriculum Development Context}
A curriculum would have to be built from the historical wordlists and the Duungidjawu grammar, working with the descendant community. Sources disagree on the language's location: the State Library notes place it in the Burnett region (Cherbourg, Kingaroy, Gayndah), while Wikipedia places it near Brisbane. These claims need checking with community and archival experts. The material found is thin, and a full K-12 curriculum cannot be supported from it alone.
\end{document}